
%=====================================================================

\section{Introduction}
%=====================================================================

The Gimnasio de Pesas at the University of Puerto Rico at Mayagüez (UPRM), located at the Natatorio Complex, relies on manual registration and access-verification steps and uses paper-based records for registration. Staff must confirm whether a person is authorized to enter through a paid membership or athlete authorization. Because these authorization sources come from separate processes, users may need to provide proof of payment while staff check different records before granting access. The records also do not automatically provide an active-user count or facility-use metrics such as peak usage periods, occupancy patterns, or differences among user groups.

A stakeholder assessment on September 11, 2026, with the Director of UPRM's Department of Athletic Activities, José F. Estévez Pérez, and two Athletic Activities staff members who work alongside the Director and are familiar with the current operational processes, confirmed these concerns. They reported reliability problems with manual access verification, said that facility-use metrics matter for administrative decisions and for students, and said users need to hear about schedule changes, temporary closures, maintenance, and other operational updates. The Director estimated that about 150 people use the Weight Gym each day, so staff repeat these verification and registration steps about 150 times a day. At the same time, some information needed for access verification already exists digitally. UPRM uses WooCommerce to manage gym membership payments, and a completed order indicates a successful payment. However, staff still perform payment verification and check-in manually, while the process does not automatically use check-in records to calculate occupancy or generate utilization metrics.

Since the proposal, the team has worked with two more stakeholders. On October 22, 2026, three team members met the contact at UPRM's Centro de Tecnologías de Información (CTI), the office that will host the application. CTI said it uses Laravel and suggests it to other projects, left the database and containers to the team, and said that the project owner, the gym director, files the virtual machine, domain, and VPN requests and answers for the accounts he authorizes. CTI also said it has no staff to maintain the application after December, so the owner must plan who looks after it. The minutes are in the project repository (\texttt{docs/meetings/2026-10-22-cti-minuta.md}). On November 4, 2026, the team meets the gym director to settle the open questions about the athlete roster, Check In rules, and staff accounts: \pending{director meeting outcomes}. This report shows what those meetings changed, what the team has built, and what remains before the final report on December 11.


%=====================================================================
